\documentclass[a4paper,11pt]{scrartcl}
\usepackage[T1]{fontenc}       % modern font encoding
\usepackage[utf8]{inputenc}    % UTF-8 encoding
\usepackage[english]{babel}    % change to german if needed
\usepackage{lmodern}           % scalable, modern font
\usepackage{graphicx}          % include figures
\usepackage{caption}           % modern caption control
\usepackage{microtype}         % better spacing
\usepackage{float}
\usepackage[
  headsepline=0.5pt,
  footsepline=0.5pt
]{scrlayer-scrpage}

% First run R_16S_AW_Butterfly_pipeline_v01.R to create figures, then run Supplementary_Information.tex 

\usepackage[
paper=a4paper,
left=15mm,
right=15mm,
top=20mm,
bottom=30mm
]{geometry}

\graphicspath{{plots_peru/}}
\renewcommand{\familydefault}{\sfdefault}
\renewcommand{\thefigure}{S\arabic{figure}} % S1, S2...
\addto\captionsenglish{%
	\renewcommand{\figurename}{Supplementary Figure}%
}
\captionsetup{%font={small},
	 labelfont={bf}, format=plain }

\ihead{Supplementary Information}
\ifoot{Weinhold \textit{et al.}, 2026}
\cfoot{}
\ofoot{\pagemark}
\setkomafont{pagehead}{\normalfont}
\setkomafont{pagefoot}{\normalfont}

\usepackage[pdfauthor={Weinhold et al. 2026},
pdftitle={Supplementary Information},
pdfsubject={Supplementary Information},
pdfpagelayout=SinglePage
]{hyperref}

\begin{document}

\begin{titlepage}
	\centering
	{\bfseries\LARGE Host-structured adult butterfly microbiomes share a common core across continents\par}
	\vspace{1.5cm}
	{\large Weinhold A$^{1*}$, Pinos A$^{1,2\#}$, Correa-Carmona Y$^{3\#}$, Holzmann KL$^{4}$, Alonso-Alonso P$^{4}$, Yon F$^{5}$, Steffan-Dewenter I$^{4}$, Peters MK$^{4,6}$, Brehm G$^{3}$, Keller A$^{1}$\par}
	\vspace{0.5cm}
	{\small
		$^{1}$Center for Organismic Adaptation, Cellular and Organismic Networks, Ludwig-Maximilians-Universität München, Munich, Germany\par
		$^{2}$Institute of Clinical Chemistry and Laboratory Medicine, University Hospital Carl Gustav Carus Dresden, Dresden, Germany\par
		$^{3}$Institut für Zoologie und Evolutionsbiologie mit Phyletischem Museum, Friedrich-Schiller University Jena, Jena, Germany\par
		$^{4}$Department of Animal Ecology and Tropical Biology, Biocenter, University of Würzburg, Würzburg, Germany\par
        $^{5}$Departamento de Ciencias Biológicas y Fisiológicas, Facultad de Ciencias e Ingeniería, Universidad Peruana Cayetano Heredia, Lima, Peru\par
        $^{6}$Animal Ecology, BIOM, University of Bremen, Germany\par
	}
    \vspace{0.5cm}
    {\small
		$^{*}$Corresponding author: arne.weinhold@bio.lmu.de\par
        $^{\#}$authors contributed equally\par
        }
    \vspace{2cm}
    Supplementary Figures S1--S11 \\
    Supplementary Table S1
    \vfill
	{\Large \textbf{Supplementary Information}\par}
	\vspace{1cm}
	{\small \today\par} % optional date
\end{titlepage}


\section*{Supplementary Figures}
% -----------------------------
% Figure S1
\pdfbookmark[1]{Figure S1 -- Bacterial load and diversity across butterfly host taxa and countries}{fig:S1}
\noindent
\begin{figure}[H]
	\includegraphics[width=\textwidth,height=0.85\textheight,keepaspectratio]{FigS1_alpha.country.pdf} 
	\caption{\textbf{Bacterial load and diversity across butterfly host taxa and countries.}
		\newline
		(\textbf{a}) Total bacterial load estimated by 16S rRNA gene copy number qPCR per butterfly genus. (\textbf{b}) Bacterial load did not differ between countries in individual subfamilies. Wilcoxon rank-sum tests: Coliadinae ($p = 0.051$), Dismorphiinae ($p = 0.33$) and Pierinae ($p = 0.81$). 
		(\textbf{c}) Shannon Index grouped by host tribe and colored by host subfamily. The full model explained 30\%{} of variation in Shannon diversity, with the main effect accounted to host tribe
		($R^2_{\mathrm{partial}} = 0.24,\; F_{11,154} = 4.3,\; p < 10^{-4}$), there was no effect of country ($R^2_{\mathrm{partial}} < 0.01,\; F_{1,154} = 0.03,\; p = 0.87$). 
		(\textbf{d}) Shannon Index did not differ by country in individual Wilcoxon rank-sum tests: Coliadinae ($p = 0.12$), Dismorphiinae ($p = 0.33$) and Pierinae ($p = 0.62$).}
\end{figure} 
\clearpage
% -----------------------------
% Figure S2
\pdfbookmark[1]{Figure S2 -- Influence of environmental gradients on adult butterfly microbiota diversity}{fig:S2}
\noindent
\begin{figure}[H]
\includegraphics[width=\textwidth,height=0.85\textheight,keepaspectratio]{FigS2_alphaelevation.pdf} 
\caption{\textbf{Influence of environmental gradients on adult butterfly microbiota diversity.}
	\newline
(\textbf{a}) Alpha diversity (only Peru samples) shows linear decrease with elevation (meters above sea level). 
(\textbf{b}) Alpha diversity (only Germany samples) shows linear increase by calendar week.
\textit{Aglais} were removed from diversity analyses, due to their low microbial abundance, but illustrated as outgroup with in the figure. 
(\textbf{c}) Temperature did not significantly affect bacterial load after accounting for host subfamily. 
(\textbf{d}) For observed microbial richness, both host subfamily and temperature were significant, with host subfamily having the stronger effect ($R^2_{\mathrm{partial}} = 0.27, F_{6,151} = 9.42, p < 10^{-8}$), while temperature had a smaller but significant effect ($R^2_{\mathrm{partial}} = 0.03, F_{1,151} = 4.63, p = 0.033$).
(\textbf{e}) For Shannon diversity, the effect was attributable to host subfamily, whereas temperature was not significant. 
(\textbf{f}) For inverse Simpson diversity, host subfamily had the stronger effect ($R^2_{\mathrm{partial}} = 0.14, F_{6,151} = 3.97, p = 0.001$), while temperature was still significant ($R^2_{\mathrm{partial}} = 0.03, F_{1,151} = 4.32, p = 0.039$).}
\end{figure} 
\clearpage
% -----------------------------
% Figure S3
\pdfbookmark[1]{Figure S3 -- Adult butterfly microbiome dissimilarity and dispersion across host genera}{fig:S3}
\noindent
\begin{figure}[H]
\includegraphics[width=\textwidth,height=0.85\textheight,keepaspectratio]{FigS3_beta_genus.pdf}
\caption{\textbf{Adult butterfly microbiome dissimilarity and dispersion across host genera.}
	\newline
(\textbf{a}) Beta diversity of the adult butterfly microbiota using PCoA based on Bray-Curtis distances. Samples are colored by host tribe and panels faceted by host genus (unique genera with less than two replicates were excluded).
(\textbf{b}) Microbiota of adult butterfly genera differs in beta dispersion (colored by host tribe).}
\end{figure} 
\clearpage
% -----------------------------
% Figure S4
\pdfbookmark[1]{Figure S4 -- Selection of the common genus core microbiota }{fig:S4}
\noindent
\begin{figure}[H]
\includegraphics[width=\textwidth,height=0.85\textheight,keepaspectratio]{FigS4_corevenn_alternative.pdf}
\caption{\textbf{Selection of the common genus core microbiota.}
	\newline
(\textbf{a}) Heatmap highlighting the most abundant and prevalent genera among all butterfly samples. (\textbf{b}) Venn diagram shows overlap of common core genera when applying different core definitions. High abundance definition  (5\% rel. abundance 10\% prevalence) results in 11 core genera (missing \textit{Acinetobacter}); Medium abundance definition (1\% rel. abundance 20\% prevalence) results in 12 core genera (as used in this study). Low abundance definition (0.2\% rel. abundance 40\% prevalence) results in 12 core genera replacing \textit{Entomomonas} and \textit{Apibacter} with \textit{Carnobacterium} and \textit{Raoultella}. 
(\textbf{c}) Average relative abundance of the 12 common genus core taxa among individual butterfly genera. (\textbf{d}) Remaining non-core microbiota among individual butterfly genera. Even non-core genera belonged to the same bacterial groups e.g. LAB (\textit{Carnobacterium}), enteric bacteria (\textit{Cedecea}, \textit{Serratia}, \textit{Roultella}), or Orbacea (\textit{Frischella}).}
\end{figure} 
\clearpage
% -----------------------------
% Figure S5
\pdfbookmark[1]{Figure S5 -- Adult butterfly microbiome composition across individual samples}{fig:S5}
\noindent
\begin{figure}[H]
	\includegraphics[width=\textwidth,height=0.95\textheight,keepaspectratio]{FigS5_sample_core_noncore.pdf}
\caption{\textbf{Adult butterfly microbiome composition across individual samples.}
	\newline
(\textbf{a}) Average relative abundance of the 12 common genus core taxa per individual sample grouped by host subfamily.
(\textbf{b}) Remaining non-core genera per sample grouped by host subfamily. Relative abundances <1\%{} are not shown.}
\end{figure} 
\clearpage
% -----------------------------
% Figure S6
\pdfbookmark[1]{Figure S6 -- Full genus-level annotation, prevalence and abundance of the adult butterfly microbiota}{fig:S6}
\noindent
\begin{figure}[H]
\includegraphics[width=\textwidth,height=0.85\textheight,keepaspectratio]{FigS6_core_square.pdf}
\caption{\textbf{Full genus-level annotation, prevalence and abundance of the adult butterfly microbiota.}
	\newline
Phylogenetic grouping (SINA RAxML tree) and relative abundance of the 80 most abundant bacterial genera from adult butterflies. Expansion of Figure 4 including all bacterial genus names. Prevalence in wild bee microbiomes from sympatric Apinae and Halictinae is shown for comparison. Black dots indicate common core genera identified from the bee dataset.}
\end{figure} 
\clearpage
% -----------------------------
% Figure S7
\pdfbookmark[1]{Figure S7 -- Comparison of the common genus core microbiota of Lepidoptera and Hymenoptera}{fig:S7}
\noindent
\begin{figure}[H]
	\includegraphics[width=\textwidth,height=0.85\textheight,keepaspectratio]{FigS7_Hym_Lep_core_stripchart.png}
	\caption{\textbf{Comparison of the common genus core microbiota of Lepidoptera and Hymenoptera}.
		\newline
	(\textbf{a}) Six genera are unique to the common core of Lepidoptera. 
	(\textbf{b}) Six genera show taxonomic overlap between the common cores of Lepidoptera and Hymenoptera. 
	(\textbf{c}) Nine genera are unique to the common core of Hymenoptera. Using the same operational common core definition for Lepidoptera and Hymenoptera should facilitate comparability between datasets but does not necessarily represent an optimized threshold for the Hymenoptera dataset, nor is it intended to specifically identify socially transmitted taxa. Wilcoxon rank-sum tests with Benjamini–Hochberg correction for multiple comparisons (Lepidoptera n = 167, Hymenoptera n = 872).}
\end{figure} 
\clearpage
% -----------------------------
% Figure S8
\pdfbookmark[1]{Figure S8 -- Country-specific distribution of common genus core taxa}{fig:S8}
\noindent
\begin{figure}[H]
\includegraphics[width=\textwidth,height=0.85\textheight,keepaspectratio]{FigS8_country_core.pdf}
\caption{\textbf{Country-specific distribution of common genus core taxa. }
	\newline
Relative abundance of individual core genera among countries and results of Wilcoxon rank-sum Tests following Benjamini–Hochberg correction accounting for multiple comparisons (Germany n = 76, Peru n = 91).}
\end{figure} 
\clearpage
% -----------------------------
% Figure S9
\pdfbookmark[1]{Figure S9 -- Distribution of common genus core ASVs across countries}{fig:S9}
\noindent
\begin{figure}[H]
\includegraphics[width=\textwidth,height=0.85\textheight,keepaspectratio]{FigS9_bubble.country.pdf} 
\caption{\textbf{Distribution of common genus core ASVs across countries.}
	\newline
 Relative abundance and distribution of individual ASVs of the 12 common genus core taxa. Only ASVs with $\ge 2\%$ relative abundance in at least one sample were shown.}
\end{figure} 
\clearpage
% -----------------------------
% Figure S10
\pdfbookmark[1]{Figure S10 -- Endosymbionts in the adult butterfly microbiota}{fig:S10}
\noindent
\begin{figure}[H]
\includegraphics[width=\textwidth,height=0.85\textheight,keepaspectratio]{FigS10_endo.pdf}
\caption{\textbf{Endosymbionts in the adult butterfly microbiota.}
	\newline
	(\textbf{a}) Average relative abundance of endosymbionts grouped per butterfly genus. (\textbf{b}) Relative abundance of endosymbionts among individual samples grouped by genus and country. Relative abundances <1\%{} are not shown.}
\end{figure} 
\clearpage
% -----------------------------
% Figure S11
\pdfbookmark[1]{Figure S11 -- Distribution of common genus core taxa among butterfly samples}{fig:S11}
\noindent
\begin{figure}[H]
\includegraphics[width=\textwidth,height=0.85\textheight,keepaspectratio]{FigS11_corestrip.pdf}
\caption{\textbf{Distribution of common genus core taxa among butterfly samples.}
	\newline
(\textbf{a}) Total genus core abundance among individual butterfly genera. 
(\textbf{b}) Distribution of individual genus core taxa across butterfly subfamilies. Symbols depict different sampling countries (Circles for Germany, triangles for Peru).}
\end{figure} 
\clearpage
% -----------------------------
\noindent
\pdfbookmark[1]{Table S1 -- Sample metadata}{tab:S1}
\textbf{Supplementary Table S1} (File: \texttt{Suppl\_table\_sample\_filter\_metadata.csv}) \\[2mm]
Sample metadata and removed reads per individual sample (total and percentage) including individual BioProject accession numbers.
\clearpage

\end{document}